% Lecture 8 notes: \input by notes.tex, \subimport'ed by the semester book.
% Exam questions (зачёт): 17 (простые числа, их бесконечность, теорема Вильсона и тест),
% 18-19 (псевдопростые, малая теорема Ферма и тест Ферма, числа Кармайкла, Миллер-Рабин),
% 20 (Соловей-Штрассен).
\begin{lecture}{8}{Простые числа и тесты простоты}

\section{Зачем нужны тесты простоты}

Почти каждый протокол с открытым ключом начинается с выбора \emph{большого простого числа}.
Для системы Диффи"--~Хеллмана нужно простое $p$ длиной $2048$ бит и больше, для RSA "--- два
случайных простых числа по $1024$ бита и больше, для эллиптических кривых "--- простое поле
$\GF{p}$. Таблиц таких чисел не бывает: их приходится \emph{генерировать} "--- брать случайное
число нужной длины и проверять, простое ли оно. Если нет, берётся следующее кандидатное число.

Тем самым нужны ответы на два вопроса.
\begin{enumerate}
  \item \emph{Как часто встречаются простые числа?} Иначе говоря, сколько случайных кандидатов
    придётся проверить. Ответ даёт теорема о распределении простых чисел: среди чисел порядка
    $x$ доля простых равна примерно $1/\ln x$; для $1024$-битных нечётных чисел это примерно
    $1/355$.
  \item \emph{Как быстро проверить простоту?} Пробное деление на все $d \le \sqrt N$ требует
    порядка $\sqrt N = 2^{512}$ операций "--- это невозможно. Нужны алгоритмы, полиномиальные по
    \emph{длине} числа $\log N$. Мы построим их из малой теоремы Ферма и её усилений:
    тест Ферма, тест Соловея"--~Штрассена и тест Миллера"--~Рабина. Тест Миллера"--~Рабина
    используется на практике везде "--- от OpenSSL до систем компьютерной алгебры.
\end{enumerate}

Тесты этой лекции \emph{вероятностные}. Если тест объявил число составным, это точно так, и
у теста есть «свидетель» "--- проверяемое доказательство. Если же тест его составным не
объявил, число лишь \emph{вероятно простое}: вероятность ошибки можно сделать сколь угодно
малой, например меньше $2^{-128}$, но не нулевой.

\section{Простые числа и их распределение}

\begin{definition}
Натуральное число $p > 1$ называется \emph{простым}, если у него нет натуральных делителей,
кроме $1$ и $p$. Натуральное число $n > 1$, не являющееся простым, называется \emph{составным}:
$n = ab$, $1 < a, b < n$. Число $1$ не является ни простым, ни составным.
\end{definition}

Основная теорема арифметики: каждое натуральное $n > 1$ однозначно (с точностью до порядка
сомножителей) представляется в виде произведения простых чисел, $n = p_1^{e_1} \cdots p_k^{e_k}$.
Её мы считаем известной из курса алгебры; единственность разложения следует из леммы Евклида
$p \mid ab \Rightarrow p \mid a$ или $p \mid b$ (она вытекает из соотношения Безу, см.\ лекцию~3).
Обозначим через $p_1 = 2, p_2 = 3, p_3 = 5, \ldots$ простые числа в порядке возрастания и через
$\pi(x)$ "--- количество простых чисел, не превосходящих $x$.

\subsection{Бесконечность множества простых чисел}

\begin{theorem}[Евклид]\label{l8:euclid}
Простых чисел бесконечно много.
\end{theorem}

\begin{proof}
Пусть $p_1, \ldots, p_n$ "--- любые простые числа. Число $N = p_1 p_2 \cdots p_n + 1 > 1$ имеет
хотя бы один простой делитель $q$. Если бы $q = p_i$ для некоторого $i$, то $q$ делил бы
$N - p_1 \cdots p_n = 1$, что невозможно. Значит, $q$ "--- простое число, отличное от
$p_1, \ldots, p_n$, и никакой конечный набор не исчерпывает все простые числа.
\end{proof}

Доказательство Евклида даёт и (очень слабую) количественную оценку.

\begin{corollary}\label{l8:loglog}
$p_n \le 2^{2^{n-1}}$ при всех $n \ge 1$; следовательно, $\pi(x) > \log_2 \log_2 x$ при $x \ge 2$.
\end{corollary}

\begin{proof}
Индукция по $n$: $p_1 = 2 = 2^{2^0}$. Наименьший простой делитель числа $p_1 \cdots p_n + 1$
отличен от $p_1, \ldots, p_n$, поэтому
\[
  p_{n+1} \le p_1 p_2 \cdots p_n + 1 \le 2^{2^0 + 2^1 + \cdots + 2^{n-1}} + 1 = 2^{2^n - 1} + 1
  \le 2^{2^n}.
\]
Пусть $x \ge 2$ и $n = \floor{\log_2 \log_2 x} + 1 \ge 1$. Тогда
$p_n \le 2^{2^{n-1}} \le 2^{\log_2 x} = x$, т.\,е. не превосходящих $x$ простых не меньше $n$, и
$\pi(x) \ge n > \log_2 \log_2 x$.
\end{proof}

\checkyourself{Почему число $p_1 p_2 \cdots p_n + 1$ из доказательства теоремы Евклида
не обязано быть простым? Приведите пример. Докажите аналогично, что простых чисел вида
$4k + 3$ бесконечно много (рассмотрите $4p_1 \cdots p_n - 1$).}

\subsection{Теорема Чебышёва и асимптотический закон}

Оценка $\log_2 \log_2 x$ безнадёжно мала. Настоящий порядок роста $\pi(x)$ "--- это
$x / \ln x$; верхняя и нижняя оценки этого порядка с явными константами получаются элементарно
из свойств центрального биномиального коэффициента $\binom{2n}{n}$.

\begin{lemma}\label{l8:binom}
Пусть $n \ge 1$. Тогда
\begin{enumerate}
  \item $2^n \le \binom{2n}{n} \le 4^n$;
  \item если $p^{e}$ "--- наибольшая степень простого $p$, делящая $\binom{2n}{n}$, то
    $p^{e} \le 2n$; в частности, $\binom{2n}{n}$ делится только на простые $p \le 2n$;
  \item каждое простое $p$, $n < p \le 2n$, делит $\binom{2n}{n}$.
\end{enumerate}
\end{lemma}

\begin{proof}
1. $\binom{2n}{n} = \prod_{k=1}^{n} \frac{n + k}{k}$ и каждый множитель не меньше $2$;
$\binom{2n}{n}$ "--- одно из слагаемых суммы $\sum_k \binom{2n}{k} = 2^{2n}$.

2. По формуле Лежандра показатель степени $p$ в $m!$ равен $\sum_{j \ge 1} \floor{m / p^j}$,
поэтому
\[
  e = \sum_{j \ge 1} \Bigl(\floor{\tfrac{2n}{p^j}} - 2\floor{\tfrac{n}{p^j}}\Bigr).
\]
Для любого вещественного $y$ разность $\floor{2y} - 2\floor{y}$ равна $0$ или $1$, а при
$p^j > 2n$ слагаемое равно нулю. Значит, $e$ не превосходит числа тех $j$, для которых
$p^j \le 2n$, т.\,е. $p^e \le 2n$.

3. Простое $p \in (n, 2n]$ входит в $(2n)!$ (как сомножитель $p$), но не входит в $(n!)^2$,
так как все сомножители $n!$ меньше $p$.
\end{proof}

\begin{theorem}[Чебышёв, 1850]\label{l8:chebyshev}
При всех $x \ge 2$
\[
  \frac{\ln 2}{2} \cdot \frac{x - 1}{\ln x} \;\le\; \pi(x) \;\le\; 8\,\frac{x}{\ln x}.
\]
\end{theorem}

\begin{proof}
\emph{Нижняя оценка.} Пусть $n = \floor{x/2} \ge 1$, так что $n \ge (x - 1)/2$ и $2n \le x$.
По лемме~\ref{l8:binom} $\binom{2n}{n} = \prod_{p \le 2n} p^{e_p}$ и каждый множитель
$p^{e_p} \le 2n$, поэтому
\[
  2^n \le \binom{2n}{n} \le (2n)^{\pi(2n)}, \qquad\text{откуда}\qquad
  \pi(x) \ge \pi(2n) \ge \frac{n \ln 2}{\ln (2n)} \ge \frac{\ln 2}{2} \cdot \frac{x - 1}{\ln x}.
\]

\emph{Верхняя оценка.} Положим $\vartheta(x) = \sum_{p \le x} \ln p$. По пп.~1 и~3 леммы
$\prod_{n < p \le 2n} p \le \binom{2n}{n} \le 4^n$, т.\,е. $\vartheta(2n) - \vartheta(n) \le n \ln 4$.
Складывая эти неравенства для $n = 1, 2, 4, \ldots, 2^{k-1}$, получаем
$\vartheta(2^k) \le (2^k - 1) \ln 4 < 2^k \ln 4$. Для $x \ge 1$ выберем $k$ с
$2^{k-1} \le x < 2^k$: $\vartheta(x) \le \vartheta(2^k) < 2^k \ln 4 \le 2x \ln 4$.
Каждое простое $p \in (\sqrt x, x]$ даёт в $\vartheta(x)$ слагаемое $\ln p > \frac12 \ln x$,
поэтому
\[
  \pi(x) \le \sqrt x + \frac{\vartheta(x)}{\frac12 \ln x} < \sqrt x + \frac{4 \ln 4 \cdot x}{\ln x}
  \le (2 + 4\ln 4)\,\frac{x}{\ln x} < 8\,\frac{x}{\ln x};
\]
здесь использовано $\sqrt x \le 2x / \ln x$, т.\,е. $\ln x \le 2\sqrt x$, и $2 + 4\ln 4 \approx 7{,}55$.
\end{proof}

Точная асимптотика была угадана Гауссом и Лежандром по таблицам простых чисел, а доказана лишь
в 1896~г. методами комплексного анализа.

\begin{theorem}[асимптотический закон распределения простых чисел; Адамар, Валле"--~Пуссен]
\label{l8:pnt}
$\displaystyle \pi(x) \sim \frac{x}{\ln x}$ при $x \to \infty$. \textup{(Без доказательства.)}
\end{theorem}

\begin{center}
\begin{tabular}{r|r|r|c}
  $x$ & $\pi(x)$ & $x / \ln x$ & $\pi(x) \big/ \frac{x}{\ln x}$ \\ \hline
  $10^3$ & $168$ & $144{,}8$ & $1{,}160$ \\
  $10^6$ & $78\,498$ & $72\,382{,}4$ & $1{,}084$ \\
  $10^9$ & $50\,847\,534$ & $48\,254\,942{,}4$ & $1{,}054$
\end{tabular}
\end{center}

Сходимость медленная; лучшее приближение даёт интегральный логарифм
$\mathrm{li}(x) = \int_2^x \frac{dt}{\ln t}$ (например, $\mathrm{li}(10^9) - \pi(10^9) \approx 1{,}7 \cdot 10^3$).
Для криптографии важно следствие: \emph{вероятность того, что случайное число около $x$ простое,
примерно равна $1/\ln x$}, а для случайного нечётного числа "--- $2/\ln x$. Для $x = 2^{1024}$
имеем $\ln x \approx 709{,}8$, т.\,е. простым окажется в среднем одно из $\approx 355$
случайных нечётных $1024$-битных чисел. Генерация простого числа нужной длины требует
лишь нескольких сотен проверок простоты.

\checkyourself{Оцените с помощью теоремы~\ref{l8:pnt}, сколько существует $512$-битных простых
чисел. Можно ли составить их таблицу? Можно ли надеяться найти «секретное» простое число
RSA перебором по такой таблице?}

\section{Теорема Вильсона}

Начнём с критерия простоты, который, в отличие от последующих тестов, является точным
(необходимым и достаточным условием), "--- но, как мы увидим, бесполезен для больших чисел.

\begin{lemma}\label{l8:sqrt1}
Если $p$ простое и $\congr{x^2}{1}{p}$, то $\congr{x}{\pm 1}{p}$.
\end{lemma}

\begin{proof}
$p \mid x^2 - 1 = (x - 1)(x + 1)$, и по лемме Евклида $p$ делит $x - 1$ или $x + 1$.
\end{proof}

Для составного модуля лемма неверна: $\congr{4^2}{1}{15}$, $\congr{11^2}{1}{15}$. Именно на
\emph{нетривиальных квадратных корнях из единицы} будет построен тест Миллера"--~Рабина.

\begin{theorem}[Вильсон, Лагранж]\label{l8:wilson}
Натуральное $n > 1$ простое тогда и только тогда, когда $\congr{(n - 1)!}{-1}{n}$.
\end{theorem}

\begin{proof}
\emph{Необходимость.} Пусть $n = p$ простое; при $p = 2$ утверждение проверяется
непосредственно: $1! = 1 \equiv -1 \pmod 2$. Пусть $p > 2$. Каждый элемент
$a \in \{1, \ldots, p - 1\}$ обратим по модулю $p$, причём $a^{-1} = a$ лишь при
$\congr{a^2}{1}{p}$, т.\,е. по лемме~\ref{l8:sqrt1} при $a = 1$ и $a = p - 1$. Остальные $p - 3$
чисел $2, \ldots, p - 2$ разбиваются на пары $\{a, a^{-1}\}$, $a \ne a^{-1}$, произведение
в каждой паре сравнимо с $1$. Поэтому
\[
  (p - 1)! \equiv 1 \cdot \underbrace{\prod_{\text{пары}} a \cdot a^{-1}}_{\equiv 1} \cdot (p - 1)
  \equiv -1 \pmod p.
\]

\emph{Достаточность.} Пусть $n$ составное, $n = ab$, $1 < a \le b < n$. Если $\congr{(n - 1)!}{-1}{n}$,
то и $\congr{(n - 1)!}{-1}{a}$; но $a \le n - 1$ входит в произведение $(n - 1)!$, т.\,е.
$\congr{(n - 1)!}{0}{a}$, и $a \mid 1$ "--- противоречие.
\end{proof}

Для составного $n$ можно сказать больше: \emph{$(n - 1)! \equiv 0 \pmod n$ при всех составных
$n > 4$}. Действительно, если $n = ab$ с $1 < a < b < n$, то $a$ и $b$ "--- различные
сомножители $(n - 1)!$. Если же $n = p^2$ и $p > 2$, то $p < 2p < p^2 = n$, и в $(n - 1)!$
входят оба сомножителя $p$ и $2p$. Единственное исключение "--- $n = 4$: $3! = 6 \equiv 2 \pmod 4$.

\begin{algorithm}[тест Вильсона]
\begin{algorithmic}[1]
\Require целое $n > 1$.
\Ensure «простое» или «составное».
\State $f \gets 1$
\For{$k = 2, 3, \ldots, n - 1$}
  \State $f \gets f \cdot k \bmod n$
  \If{$f = 0$} \Return «составное» \EndIf
\EndFor
\If{$f = n - 1$} \Return «простое» \Else\ \Return «составное» \EndIf
\end{algorithmic}
\end{algorithm}

\taskexample{wilson-test}

Тест Вильсона \emph{детерминированный и точный}, но для простого $n$ он выполняет $n - 2$
умножений по модулю $n$, т.\,е. $\Theta(n \log^2 n)$ битовых операций. Это экспонента от
длины входа $\log_2 n$: для $n \approx 2^{1024}$ "--- около $2^{1024}$ умножений, что намного
хуже даже пробного деления ($\sqrt n$ делений). Быстрого способа вычислять $(n - 1)! \bmod n$
неизвестно: лучшие алгоритмы требуют порядка $\sqrt n$ операций (с логарифмическими
множителями), что не лучше пробного деления. Теорема Вильсона "--- красивый
\emph{критерий} простоты, но не практический \emph{тест}.

\checkyourself{Вычислите $(n - 1)! \bmod n$ для $n = 6, 8, 9, 25$. На каком шаге алгоритма
результат обращается в нуль? Почему для $n = 4$ этого не происходит?}

\section{Малая теорема Ферма и тест Ферма}

\begin{theorem}[Эйлер]\label{l8:euler}
Если $\gcd(a, n) = 1$, то $\congr{a^{\varphi(n)}}{1}{n}$, где $\varphi(n) = \abs{\Znx{n}}$ "---
функция Эйлера.
\end{theorem}

\begin{proof}
Пусть $r_1, \ldots, r_m$, $m = \varphi(n)$, "--- все вычеты из $\Znx{n}$. Так как $a$ обратим,
отображение $x \mapsto ax$ инъективно на $\Znx{n}$ (из $ax \equiv ay$ следует $x \equiv y$)
и, следовательно, биективно: $ar_1, \ldots, ar_m$ "--- те же вычеты в другом порядке. Поэтому
\[
  \prod_{i=1}^{m} (a r_i) \equiv \prod_{i=1}^{m} r_i \pmod n
  \quad\Longrightarrow\quad a^{m} \prod_i r_i \equiv \prod_i r_i \pmod n,
\]
и после сокращения на обратимый элемент $\prod_i r_i$ получаем $\congr{a^m}{1}{n}$.
\end{proof}

\begin{corollary}[малая теорема Ферма]\label{l8:fermat}
Если $p$ простое и $p \nmid a$, то $\congr{a^{p-1}}{1}{p}$. Для любого целого $a$:
$\congr{a^p}{a}{p}$.
\end{corollary}

\begin{proof}
$\varphi(p) = p - 1$. Второе утверждение при $p \nmid a$ получается умножением на $a$,
а при $p \mid a$ обе части сравнимы с нулём.
\end{proof}

Обратное неверно: из $\congr{a^{n-1}}{1}{n}$ простота $n$ не следует. Но по контрапозиции
\emph{если $a^{n-1} \not\equiv 1 \pmod n$, то $n$ заведомо составное} "--- и для проверки
достаточно одного возведения в степень, $\bigO(\log n)$ умножений по модулю $n$ (лекция~6).

\begin{definition}
Пусть $n > 1$ нечётное и $1 \le a < n$.
\begin{itemize}
  \item Если $a^{n-1} \not\equiv 1 \pmod n$, то $a$ называется \emph{свидетелем (Ферма)
    составности} $n$.
  \item Если $\congr{a^{n-1}}{1}{n}$, то $n$ называется \emph{вероятно простым по основанию $a$}
    (англ. probable prime); если при этом $n$ составное, то $n$ называется \emph{псевдопростым
    по основанию $a$}, а $a$ "--- \emph{лжесвидетелем} (англ. Fermat liar).
\end{itemize}
\end{definition}

Заметим, что основание $a$ с $\gcd(a, n) = d > 1$ всегда свидетель: из $\congr{a^{n-1}}{1}{n}$
следовало бы $d \mid 1$. Таким образом, лжесвидетели лежат в $\Znx{n}$.

\begin{example}\label{l8:ex341}
$n = 341 = 11 \cdot 31$ "--- наименьшее псевдопростое по основанию $2$. Действительно,
$2^{10} = 1024 = 3 \cdot 341 + 1$, т.\,е. $\congr{2^{10}}{1}{341}$ и
$2^{340} = (2^{10})^{34} \equiv 1$. Но $\congr{3^{340}}{56}{341}$, и основание $3$ "--- свидетель
составности.
\end{example}

\begin{algorithm}[тест Ферма]\label{l8:alg-fermat}
\begin{algorithmic}[1]
\Require нечётное $n > 3$, число раундов $k$.
\Ensure «составное» (наверняка) или «вероятно простое».
\For{$i = 1, \ldots, k$}
  \State выбрать случайное $a \in \{2, \ldots, n - 2\}$
  \If{$\gcd(a, n) > 1$ или $a^{n-1} \bmod n \ne 1$} \Return «составное» \EndIf
\EndFor
\State \Return «вероятно простое»
\end{algorithmic}
\end{algorithm}

При школьном умножении один раунд стоит $\bigO(\log^3 n)$ битовых операций, весь тест "---
$\bigO(k \log^3 n)$. Насколько надёжен ответ «вероятно простое»? Ключевое наблюдение:
лжесвидетели образуют подгруппу.

\begin{proposition}\label{l8:fermat-half}
Множество $L_F(n) = \{a \in \Znx{n} : a^{n-1} \equiv 1\}$ "--- подгруппа в $\Znx{n}$. Если
существует хотя бы один свидетель Ферма $a \in \Znx{n}$, то $\abs{L_F(n)} \le \varphi(n)/2$.
\end{proposition}

\begin{proof}
$1 \in L_F(n)$, и если $a^{n-1} \equiv b^{n-1} \equiv 1$, то $(ab^{-1})^{n-1} \equiv 1$.
Если $L_F(n) \ne \Znx{n}$, то по теореме Лагранжа $\abs{L_F(n)}$ "--- собственный делитель
$\varphi(n)$, т.\,е. не больше $\varphi(n)/2$.
\end{proof}

Для свободных от квадратов $n$ число лжесвидетелей можно вычислить точно.

\begin{proposition}\label{l8:fermat-count}
Если $n = p_1 p_2 \cdots p_k$ "--- произведение различных нечётных простых, то
\[
  \abs{L_F(n)} = \prod_{i=1}^{k} \gcd(n - 1, p_i - 1).
\]
\end{proposition}

\begin{proof}
По китайской теореме об остатках (лекция~5) $\Znx{n} \cong \Znx{p_1} \times \cdots \times \Znx{p_k}$,
и $a^{n-1} \equiv 1 \pmod n$ тогда и только тогда, когда $a^{n-1} \equiv 1 \pmod {p_i}$ для всех
$i$. Группа $\Znx{p}$ циклическая порядка $p - 1$ (существует первообразный корень $g$,
лекция~7), и $g^{t\,(n-1)} = 1$ $\iff$ $p - 1 \mid t(n - 1)$ $\iff$
$\frac{p - 1}{\gcd(n - 1, p - 1)} \mid t$. Таких $t \in \{0, \ldots, p - 2\}$ ровно
$\gcd(n - 1, p - 1)$. Перемножая числа решений по всем $p_i$, получаем утверждение.
\end{proof}

Например, для $n = 341$: $\abs{L_F} = \gcd(340, 10) \cdot \gcd(340, 30) = 10 \cdot 10 = 100$
из $\varphi(341) = 300$, т.\,е. одна треть.

Итак, если у $n$ есть хотя бы один свидетель, взаимно простой с $n$, то случайное основание
оказывается свидетелем с вероятностью не меньше $1/2$, и $k$ раундов ошибаются с вероятностью
не более $2^{-k}$. Но бывают составные $n$, у которых \emph{вообще нет} взаимно простых
свидетелей.

\checkyourself{Почему в определении свидетеля Ферма основания $a = 1$ и $a = n - 1$ не
интересны? Чему равно $(n - 1)^{n-1} \bmod n$ для нечётного $n$?}

\section{Числа Кармайкла}

\begin{definition}
Составное число $n$ называется \emph{числом Кармайкла}, если $\congr{a^{n-1}}{1}{n}$ для всех
$a$, взаимно простых с $n$, т.\,е. $L_F(n) = \Znx{n}$.
\end{definition}

Наименьшее число Кармайкла "--- $561 = 3 \cdot 11 \cdot 17$; числа Кармайкла, меньшие $10^4$:
$561, 1105, 1729, 2465, 2821, 6601, 8911$. Для них тест Ферма ошибается при \emph{любом}
основании, взаимно простом с $n$; составность обнаруживает только основание, имеющее с $n$
общий делитель, а их доля $1 - \varphi(n)/n = 1 - \prod (1 - 1/p_i)$ ничтожна, если все
простые делители $n$ велики. Проверять, является ли $n$ числом Кармайкла, удобно по критерию
Корсельта.

\begin{theorem}[критерий Корсельта, 1899]\label{l8:korselt}
Составное $n > 2$ является числом Кармайкла тогда и только тогда, когда
\begin{enumerate}
  \item $n$ свободно от квадратов: $p^2 \nmid n$ ни для какого простого $p$, и
  \item $p - 1 \mid n - 1$ для каждого простого делителя $p$ числа $n$.
\end{enumerate}
\end{theorem}

\begin{proof}
\emph{Достаточность.} Пусть $n = p_1 \cdots p_k$ (различные простые) и $p_i - 1 \mid n - 1$.
Возьмём $a$, $\gcd(a, n) = 1$. Тогда $p_i \nmid a$, и по малой теореме Ферма
$a^{p_i - 1} \equiv 1 \pmod{p_i}$, а значит, и $a^{n-1} = (a^{p_i - 1})^{(n-1)/(p_i - 1)} \equiv 1
\pmod{p_i}$ при всех $i$. Так как $p_i$ попарно взаимно просты, $a^{n-1} \equiv 1 \pmod n$.

\emph{Необходимость.} Пусть $n$ "--- число Кармайкла.

(1) Предположим, что $p^2 \mid n$. Положим $a = 1 + n/p$. Каждый простой делитель $q$ числа $n$
делит $n/p$ (при $q = p$ "--- потому что $p^2 \mid n$), поэтому $a \equiv 1 \pmod q$ и
$\gcd(a, n) = 1$. Далее, $(n/p)^2 = n \cdot (n/p^2) \equiv 0 \pmod n$, и по формуле бинома
$a^{m} \equiv 1 + m\,\frac{n}{p} \pmod n$ для всех $m \ge 0$. Поэтому
\[
  a^{n-1} \equiv 1 + (n - 1)\frac{n}{p} \pmod n, \qquad\text{и}\qquad a^{n-1} \equiv 1
  \iff n \bigm| (n - 1)\frac{n}{p} \iff p \mid n - 1,
\]
что невозможно, так как $p \mid n$. Значит, $a$ "--- свидетель, взаимно простой с $n$, "---
противоречие.

(2) Итак, $n$ свободно от квадратов. Пусть $p \mid n$; тогда $\gcd(p, n/p) = 1$. Пусть $g$ "---
первообразный корень по модулю $p$ (лекция~7). По китайской теореме об остатках найдётся $a$
с $a \equiv g \pmod p$ и $a \equiv 1 \pmod{n/p}$; оно взаимно просто с $n$. Тогда
$a^{n-1} \equiv 1 \pmod n$, в частности $g^{n-1} \equiv 1 \pmod p$, а так как порядок $g$
равен $p - 1$, то $p - 1 \mid n - 1$.
\end{proof}

\begin{corollary}\label{l8:carm-odd}
Каждое число Кармайкла нечётно и имеет не менее трёх простых делителей.
\end{corollary}

\begin{proof}
Пусть $n$ "--- число Кармайкла. Оно свободно от квадратов и составное, поэтому имеет нечётный
простой делитель $p$; число $p - 1$ чётно и делит $n - 1$, значит, $n$ нечётно. Если бы
$n = pq$, $p < q$, то $q - 1 \mid n - 1 = p(q - 1) + (p - 1)$, откуда $q - 1 \mid p - 1$, но
$0 < p - 1 < q - 1$ "--- противоречие.
\end{proof}

\begin{example}
$561 = 3 \cdot 11 \cdot 17$, $560 = 2^4 \cdot 5 \cdot 7$ делится на $2$, $10$ и $16$ "---
число Кармайкла. Для $1729 = 7 \cdot 13 \cdot 19$ число $1728 = 2^6 \cdot 3^3$ делится на
$6$, $12$ и $18$.
\end{example}

\begin{proposition}[Чернек, 1939]\label{l8:chernick}
Если числа $6k + 1$, $12k + 1$, $18k + 1$ простые, то их произведение "--- число Кармайкла.
\end{proposition}

\begin{proof}
$n = (6k + 1)(12k + 1)(18k + 1) = 1296k^3 + 396k^2 + 36k + 1$, и $n - 1 = 36k(36k^2 + 11k + 1)$
делится на $6k$, $12k$ и $18k$. Свободность от квадратов очевидна; применяем критерий
Корсельта.
\end{proof}

При $k = 1$ получаем $1729 = 7 \cdot 13 \cdot 19$, при $k = 6$ "--- $294\,409 = 37 \cdot 73 \cdot 109$.
Бесконечно ли много таких $k$, неизвестно. Тем не менее справедлива

\begin{theorem}[Алфорд, Грэнвилл, Померанс, 1994]
Чисел Кармайкла бесконечно много; более того, $C(x) > x^{2/7}$ для всех достаточно больших
$x$, где $C(x)$ "--- количество чисел Кармайкла, не превосходящих $x$. \textup{(Без
доказательства, см.~\cite{alford-granville-pomerance}.)}
\end{theorem}

Числа Кармайкла сравнительно редки (их $16$ до $10^5$ и $43$ до $10^6$), но их бесконечно
много, и \emph{тест Ферма не может служить надёжным тестом простоты}: на числе Кармайкла
он ошибается с вероятностью, близкой к единице, сколько бы раундов ни выполнялось.

\taskexample{fermat-test-carmichael}

\taskexample{korselt-check}

\checkyourself{Почему в доказательстве необходимости в критерии Корсельта нельзя сразу
взять $a = g$ (первообразный корень по модулю $p$)? Может ли число Кармайкла делиться
на $3$ и на $5$ одновременно? на $3$ и на $7$?}

\section{Тест Миллера"--~Рабина}

Тест Ферма проверяет только последнее звено цепочки
$a^d, a^{2d}, \ldots, a^{2^s d} = a^{n-1}$. Тест Миллера"--~Рабина смотрит на всю цепочку
и проверяет, нет ли в ней нетривиального квадратного корня из единицы.

\begin{theorem}\label{l8:mr-prime}
Пусть $p$ "--- нечётное простое, $p - 1 = 2^s d$, $d$ нечётно, и $p \nmid a$. Тогда
\[
  a^d \equiv 1 \pmod p \qquad\text{или}\qquad a^{2^i d} \equiv -1 \pmod p
  \ \text{ для некоторого } 0 \le i < s.
\]
\end{theorem}

\begin{proof}
Рассмотрим последовательность $x_i = a^{2^i d} \bmod p$, $i = 0, 1, \ldots, s$; каждый
следующий член "--- квадрат предыдущего, и $x_s = a^{p-1} \equiv 1$ по малой теореме Ферма.
Если $x_0 = 1$, всё доказано. Иначе пусть $j \ge 1$ "--- наименьший индекс с $x_j = 1$; тогда
$x_{j-1} \ne 1$ и $x_{j-1}^2 = 1$, и по лемме~\ref{l8:sqrt1} $x_{j-1} \equiv -1$, где
$j - 1 < s$.
\end{proof}

\begin{definition}
Пусть $n$ нечётное, $n - 1 = 2^s d$, $d$ нечётно, $1 \le a < n$. Число $a$ называется
\emph{свидетелем Миллера} (составности $n$), если $a^d \not\equiv 1$ и
$a^{2^i d} \not\equiv -1 \pmod n$ при всех $0 \le i < s$. Если $a$ не является свидетелем,
а $n$ составное, то $n$ называется \emph{сильно псевдопростым по основанию $a$}, а $a$ "---
\emph{сильным лжесвидетелем}.
\end{definition}

По теореме~\ref{l8:mr-prime} у простого числа свидетелей нет. Каждый свидетель Ферма "---
свидетель Миллера (если $x_s = a^{n-1} \ne 1$, условие теоремы нарушено), а обратное неверно,
что и делает тест сильнее. Если $a$ "--- свидетель Миллера, но не свидетель Ферма, то
в последовательности $x_i$ найдётся $x_j \ne \pm 1$ с $x_{j+1} = 1$ "--- \emph{нетривиальный
квадратный корень из единицы}. Тогда $n \mid (x_j - 1)(x_j + 1)$, но $n$ не делит ни одну из
скобок, так что $\gcd(x_j - 1, n)$ "--- \emph{нетривиальный делитель} $n$. Тест не только
доказывает составность, но и иногда находит делитель.

\begin{algorithm}[тест Миллера"--~Рабина]\label{l8:alg-mr}
\begin{algorithmic}[1]
\Require нечётное $n > 3$, число раундов $k$.
\Ensure «составное» (наверняка) или «вероятно простое».
\State представить $n - 1 = 2^s d$, $d$ нечётно
\For{$i = 1, \ldots, k$}
  \State выбрать случайное $a \in \{2, \ldots, n - 2\}$
  \State $x \gets a^d \bmod n$
  \If{$x = 1$ или $x = n - 1$} \textbf{перейти} к следующему раунду \EndIf
  \For{$j = 1, \ldots, s - 1$}
    \State $x \gets x^2 \bmod n$
    \If{$x = n - 1$} \textbf{перейти} к следующему раунду \EndIf
  \EndFor
  \State \Return «составное» \Comment{$a$ "--- свидетель Миллера}
\EndFor
\State \Return «вероятно простое»
\end{algorithmic}
\end{algorithm}

Раунд стоит одного возведения в степень $d$ и не более $s - 1$ возведений в квадрат, т.\,е.
столько же, сколько раунд теста Ферма: $\bigO(\log^3 n)$ битовых операций при школьном
умножении.

\taskexample{miller-rabin-561}

\taskexample{miller-rabin-prime}

Главное отличие от теста Ферма: у составного числа свидетелей Миллера \emph{всегда много}.
Начнём с простой оценки, доказательство которой аналогично предложению~\ref{l8:fermat-half}.

\begin{theorem}\label{l8:mr-half}
Если $n$ нечётное составное, то число сильных лжесвидетелей не превосходит $\varphi(n)/2$.
\end{theorem}

\begin{proof}
Все сильные лжесвидетели лежат в $\Znx{n}$. Построим собственную подгруппу $B \subset \Znx{n}$,
содержащую их все.

\emph{Случай 1: $p^2 \mid n$ для некоторого простого $p$.} Сильный лжесвидетель является и
лжесвидетелем Ферма (если $x_0 = 1$ или $x_i = -1$ при $i < s$, то $x_s = 1$). В доказательстве
критерия Корсельта построен элемент $a = 1 + n/p \in \Znx{n}$, не лежащий в подгруппе $L_F(n)$.
Берём $B = L_F(n)$.

\emph{Случай 2: $n = n_1 n_2$, $\gcd(n_1, n_2) = 1$, $n_1, n_2 > 1$} (это так для любого $n$, не
являющегося степенью простого). Основание $a = -1$ даёт $(-1)^{d} = -1$, поэтому существуют
основания $a$ с $a^{2^t d} \equiv -1$; пусть $t$ "--- наибольшее такое число,
а $a_0$ "--- основание с $a_0^{2^t d} \equiv -1 \pmod n$. Положим $m = 2^t d$ и
\[
  B = \{a \in \Znx{n} : a^{m} \equiv \pm 1 \pmod n\}.
\]
Это подгруппа: она содержит $1$ и замкнута относительно умножения и взятия обратного.
Сильный лжесвидетель $a$ лежит в $B$: если $a^d \equiv 1$, то $a^{m} \equiv 1$; если
$a^{2^i d} \equiv -1$, то $i \le t$ по выбору $t$, и $a^{m} = (a^{2^i d})^{2^{t - i}} \equiv \pm 1$.
Покажем, что $B \ne \Znx{n}$. По китайской теореме найдётся $b$ с $b \equiv a_0 \pmod{n_1}$
и $b \equiv 1 \pmod{n_2}$; оно взаимно просто с $n$. Тогда $b^{m} \equiv -1 \pmod {n_1}$ и
$b^{m} \equiv 1 \pmod{n_2}$. Если бы $b^m \equiv 1 \pmod n$, то $-1 \equiv 1 \pmod{n_1}$, а
если $b^m \equiv -1 \pmod n$, то $1 \equiv -1 \pmod{n_2}$; оба варианта невозможны, так как
$n_1, n_2$ нечётны и больше $1$. Значит, $b \notin B$.

В обоих случаях $B$ "--- собственная подгруппа, и по теореме Лагранжа
$\abs{B} \le \varphi(n)/2$.
\end{proof}

На самом деле справедлива более сильная оценка.

\begin{theorem}[Монье, Рабин, 1980]\label{l8:monier-rabin}
Если $n > 9$ нечётное составное, то число сильных лжесвидетелей в $\{1, \ldots, n - 1\}$
не превосходит $\varphi(n)/4$. \textup{(Без доказательства, см.~\cite[разд.~3.5]{crandall-pomerance},
\cite{monier, rabin}.)}
\end{theorem}

\emph{Идея доказательства.} Для $n = p_1^{e_1} \cdots p_k^{e_k}$ число сильных лжесвидетелей
вычисляется явно через $\nu = \min_i \nu_2(p_i - 1)$ и $\gcd(d, p_i - 1)$ (формула Монье);
дальше разбираются случаи $k = 1$, $k = 2$ и $k \ge 3$, и в каждом показывается, что формула
даёт не больше $\varphi(n)/4$. Оценка точна: у $n = 91 = 7 \cdot 13$ ровно $18 = 72/4$ сильных
лжесвидетелей ($1, 9, 10, 12, 16, 17, 22, 29, 38, 53, 62, 69, 74, 75, 79, 81, 82, 90$).

\begin{corollary}
Если нечётное составное $n$ прошло $k$ раундов теста Миллера"--~Рабина с независимыми
случайными основаниями, вероятность этого события не превосходит $4^{-k}$.
\end{corollary}

\begin{remark}
Это оценка вероятности \emph{для фиксированного составного $n$}. Вероятность того, что
\emph{случайное} $1024$-битное число, прошедшее $k$ раундов, составное, гораздо меньше $4^{-k}$:
для случайных чисел сильные лжесвидетели встречаются крайне редко, поэтому стандарты
(например, FIPS~186) допускают всего несколько раундов для чисел такой длины.
\end{remark}

\subsection{Детерминированные варианты}

Если $n$ ограничено, достаточно перебрать фиксированный набор оснований: для каждого
набора известна наименьшая сильно псевдопростая по всем этим основаниям величина
$\psi$ (найдены компьютерным перебором; приводим без доказательства, см.~\cite{sorenson-webster}).

\begin{center}
\begin{tabular}{l|l}
  основания & тест точен для всех $n <{}$ \\ \hline
  $2$ & $2047 = 23 \cdot 89$ \\
  $2, 3$ & $1\,373\,653 = 829 \cdot 1657$ \\
  $2, 3, 5$ & $25\,326\,001 = 2251 \cdot 11\,251$ \\
  $2, 3, 5, 7$ & $3\,215\,031\,751 = 151 \cdot 751 \cdot 28\,351$ \\
  первые $12$ простых ($2, \ldots, 37$) & $318\,665\,857\,834\,031\,151\,167\,461 \approx 3{,}2 \cdot 10^{23}$ \\
  первые $13$ простых ($2, \ldots, 41$) & $3\,317\,044\,064\,679\,887\,385\,961\,981 \approx 3{,}3 \cdot 10^{24}$
\end{tabular}
\end{center}

В частности, проверка оснований $2, \ldots, 37$ даёт точный тест для всех $64$-битных чисел.
Каждое из чисел правого столбца составное и сильно псевдопростое по всем основаниям своей
строки (это легко проверить на компьютере); то, что меньших таких чисел нет, "--- результат
большого перебора.

\begin{theorem}[Миллер, 1976; Бах, 1990]
Если верна обобщённая гипотеза Римана (GRH), то у каждого нечётного составного $n$ есть
свидетель Миллера $a < 2 \ln^2 n$. \textup{(Без доказательства, см.~\cite{miller, bach-grh}.)}
\end{theorem}

Следовательно, при GRH тест Миллера с основаниями $a < 2\ln^2 n$ детерминированный и
выполняется за $\bigO(\log^5 n)$ битовых операций. Исторически Миллер предложил именно
детерминированный (условный) тест, а Рабин "--- его вероятностный вариант с безусловной
оценкой ошибки. Безусловный детерминированный полиномиальный тест простоты построили
в 2002~г. Агравал, Каял и Саксена~\cite{aks} (тест AKS, основанный на сравнении
$(x + a)^n \equiv x^n + a$ в кольце $\Zn{n}[x]/(x^r - 1)$); его современные варианты работают
за $\softO(\log^6 n)$, но на практике заметно медленнее теста Миллера"--~Рабина.

\checkyourself{Почему в тесте Миллера"--~Рабина основания $a = 1$ и $a = n - 1$ бесполезны?
Покажите, что $n = 2047$ "--- сильно псевдопростое по основанию $2$ ($2^{11} = 2048$).
Как найти делитель $n$, если известен нетривиальный квадратный корень из единицы?}

\section{Символы Лежандра и Якоби}

Ещё один тест получается из \emph{критерия Эйлера}: для простого $p$ значение $a^{(p-1)/2}$
определяется тем, является ли $a$ квадратом по модулю $p$.

\begin{definition}
Пусть $p$ "--- нечётное простое. Число $a$, $p \nmid a$, называется \emph{квадратичным вычетом}
по модулю $p$, если сравнение $x^2 \equiv a \pmod p$ разрешимо, и \emph{квадратичным невычетом}
иначе. \emph{Символ Лежандра}:
\[
  \legendre{a}{p} = \begin{cases}
    0, & p \mid a,\\
    1, & a \text{ "--- квадратичный вычет},\\
    -1, & a \text{ "--- квадратичный невычет}.
  \end{cases}
\]
\end{definition}

\begin{theorem}[критерий Эйлера]\label{l8:euler-crit}
Для нечётного простого $p$ и любого целого $a$: $\congr{a^{(p-1)/2}}{\legendre{a}{p}}{p}$.
\end{theorem}

\begin{proof}
При $p \mid a$ обе части равны нулю. Пусть $p \nmid a$ и $y = a^{(p-1)/2}$; тогда
$y^2 = a^{p-1} \equiv 1$, и по лемме~\ref{l8:sqrt1} $y \equiv \pm 1$. Если $a \equiv x^2$
квадратичный вычет, то $y \equiv x^{p-1} \equiv 1$. Квадратичных вычетов ровно $(p - 1)/2$:
отображение $x \mapsto x^2$ на $\Znx{p}$ склеивает ровно пары $\pm x$ (по лемме~\ref{l8:sqrt1},
$x^2 \equiv y^2 \iff x \equiv \pm y$). Многочлен $X^{(p-1)/2} - 1$ над полем $\GF{p}$ имеет не
более $(p - 1)/2$ корней, и все они "--- квадратичные вычеты. Значит, для невычета $a$
$y \not\equiv 1$, т.\,е. $y \equiv -1$.
\end{proof}

\begin{corollary}\label{l8:legendre-props}
\begin{enumerate}
  \item $\legendre{ab}{p} = \legendre{a}{p}\legendre{b}{p}$; $\legendre{a}{p}$ зависит только
    от $a \bmod p$.
  \item $\legendre{-1}{p} = (-1)^{(p-1)/2}$, т.\,е. $-1$ "--- квадрат по модулю $p$ $\iff$
    $p \equiv 1 \pmod 4$.
\end{enumerate}
\end{corollary}

Например, по модулю $11$ квадраты "--- это $1, 3, 4, 5, 9$, а $2^5 = 32 \equiv -1 \pmod{11}$:
$2$ "--- невычет, в согласии с критерием Эйлера.

\begin{theorem}[квадратичный закон взаимности, Гаусс]\label{l8:qr}
Для различных нечётных простых $p$ и $q$
\[
  \legendre{p}{q}\legendre{q}{p} = (-1)^{\frac{p-1}{2}\cdot\frac{q-1}{2}}, \qquad
  \legendre{2}{p} = (-1)^{\frac{p^2 - 1}{8}}.
\]
\textup{(Без доказательства; см.~\cite[разд.~3.9]{hps}, \cite[гл.~5]{ireland-rosen}.)}
\end{theorem}

Иначе говоря, $\legendre{p}{q} = -\legendre{q}{p}$, если $p \equiv q \equiv 3 \pmod 4$, и
$\legendre{p}{q} = \legendre{q}{p}$ иначе; $2$ "--- квадрат по модулю $p$ $\iff$
$p \equiv \pm 1 \pmod 8$.

\begin{definition}
Пусть $n$ "--- нечётное натуральное, $n = p_1 p_2 \cdots p_k$ (простые, возможно, с повторениями).
\emph{Символ Якоби}: $\legendre{a}{n} = \prod_{i=1}^{k} \legendre{a}{p_i}$; $\legendre{a}{1} = 1$.
\end{definition}

\begin{proposition}\label{l8:jacobi-props}
Для нечётных натуральных $m, n$:
\begin{enumerate}
  \item $\legendre{a}{n}$ зависит только от $a \bmod n$; $\legendre{ab}{n} = \legendre{a}{n}\legendre{b}{n}$,
    $\legendre{a}{mn} = \legendre{a}{m}\legendre{a}{n}$; $\legendre{a}{n} = 0 \iff \gcd(a, n) > 1$;
  \item $\legendre{-1}{n} = (-1)^{\frac{n-1}{2}}$, \quad $\legendre{2}{n} = (-1)^{\frac{n^2-1}{8}}$;
  \item если $\gcd(m, n) = 1$, то $\legendre{m}{n}\legendre{n}{m} = (-1)^{\frac{m-1}{2}\cdot\frac{n-1}{2}}$.
\end{enumerate}
\end{proposition}

\begin{proof}
1 следует из определения и свойств символа Лежандра. Для 2 и 3 заметим, что функции
$\varepsilon(n) = \frac{n-1}{2} \bmod 2$ и $\omega(n) = \frac{n^2-1}{8} \bmod 2$ аддитивны на
нечётных числах:
\begin{gather*}
  \frac{ab - 1}{2} - \frac{a - 1}{2} - \frac{b - 1}{2} = \frac{(a - 1)(b - 1)}{2} \equiv 0 \pmod 2,\\
  \frac{a^2b^2 - 1}{8} - \frac{a^2 - 1}{8} - \frac{b^2 - 1}{8} = \frac{(a^2 - 1)(b^2 - 1)}{8} \equiv 0
  \pmod 2
\end{gather*}
(числитель первой дроби делится на $4$, второй "--- на $64$). Поэтому, раскладывая $n$ на
простые и применяя следствие~\ref{l8:legendre-props} и теорему~\ref{l8:qr} к каждому
множителю, получаем 2. Аналогично, для $m = \prod q_j$, $n = \prod p_i$
\[
  \legendre{m}{n}\legendre{n}{m} = \prod_{i,j} \legendre{q_j}{p_i}\legendre{p_i}{q_j}
  = (-1)^{\sum_{i,j} \varepsilon(q_j)\varepsilon(p_i)} = (-1)^{\varepsilon(m)\varepsilon(n)}. \qedhere
\]
\end{proof}

Эти свойства дают алгоритм вычисления символа Якоби \emph{без разложения на множители},
похожий на алгоритм Евклида: выделяем из числителя множители $2$ (п.~2), переворачиваем
символ (п.~3) и приводим числитель по модулю знаменателя (п.~1).

\begin{algorithm}[вычисление символа Якоби]\label{l8:alg-jacobi}
\begin{algorithmic}[1]
\Require целое $a$, нечётное $n \ge 1$.
\Ensure $\legendre{a}{n}$.
\State $a \gets a \bmod n$, \quad $t \gets 1$
\While{$a \ne 0$}
  \While{$a$ чётно}
    \State $a \gets a/2$
    \If{$n \bmod 8 \in \{3, 5\}$} $t \gets -t$ \EndIf
  \EndWhile
  \State $(a, n) \gets (n, a)$ \Comment{взаимность}
  \If{$a \equiv n \equiv 3 \pmod 4$} $t \gets -t$ \EndIf
  \State $a \gets a \bmod n$
\EndWhile
\If{$n = 1$} \Return $t$ \Else\ \Return $0$ \EndIf
\end{algorithmic}
\end{algorithm}

Как и в алгоритме Евклида, числа убывают экспоненциально быстро; число итераций "---
$\bigO(\log n)$, сложность "--- $\bigO(\log^2 n)$ битовых операций (на каждом шаге "---
деление с остатком).

\begin{example}
$p = 9907$ простое, $1001 \equiv 1 \pmod 4$, и
\[
\begin{aligned}
  \legendre{1001}{9907} &= \legendre{9907}{1001} = \legendre{898}{1001}
  = \legendre{2}{1001}\legendre{449}{1001} = \legendre{449}{1001} = \legendre{1001}{449}
  = \legendre{103}{449} \\
  &= \legendre{449}{103} = \legendre{37}{103} = \legendre{103}{37} = \legendre{29}{37}
  = \legendre{37}{29} = \legendre{8}{29} = \legendre{2}{29}^3 = -1,
\end{aligned}
\]
так как $1001 \equiv 1 \pmod 8$, $449 \equiv 37 \equiv 29 \equiv 1 \pmod 4$ и $29 \equiv 5 \pmod 8$.
Итак, $1001$ "--- квадратичный невычет по модулю $9907$.
\end{example}

\begin{remark}
Для составного $n$ равенство $\legendre{a}{n} = 1$ \emph{не} означает, что $a$ "--- квадрат
по модулю $n$: $\legendre{2}{15} = \legendre{2}{3}\legendre{2}{5} = (-1)(-1) = 1$, но $2$
не является квадратом даже по модулю $3$.
\end{remark}

\checkyourself{Вычислите $\legendre{2}{561}$ и $\legendre{5}{21}$. Почему в алгоритме
вычисления символа Якоби не нужно раскладывать $n$ на множители, хотя символ Якоби
определён через разложение?}

\section{Тест Соловея"--~Штрассена}

По критерию Эйлера для простого $p$ и любого $a$ выполнено $a^{(p-1)/2} \equiv \legendre{a}{p}$.
Обе части вычисляются быстро и для составного $n$ (символ Якоби "--- алгоритмом выше), что
и даёт тест.

\begin{definition}
Нечётное составное $n$ называется \emph{эйлеровым псевдопростым по основанию $a$},
$\gcd(a, n) = 1$, если $a^{(n-1)/2} \equiv \legendre{a}{n} \pmod n$; $a$ называется
\emph{эйлеровым лжесвидетелем}. Иначе (в том числе при $\gcd(a, n) > 1$) $a$ "---
\emph{свидетель Эйлера}.
\end{definition}

Эйлерово псевдопростое число псевдопростое (возводим сравнение в квадрат, $\legendre{a}{n}^2 = 1$).
Например, $561$ "--- эйлерово псевдопростое по основанию $2$: $2^{280} \equiv 1 \pmod{561}$
(см.\ пример~\ref{ex:miller-rabin-561}, $x_3 = 2^{280} = 1$) и $\legendre{2}{561} = 1$, так как
$561 \equiv 1 \pmod 8$.

\begin{theorem}[Соловей, Штрассен, 1977]\label{l8:ss-half}
Если $n$ нечётное составное, то эйлеровы лжесвидетели образуют собственную подгруппу
$E(n) \subset \Znx{n}$; в частности, $\abs{E(n)} \le \varphi(n)/2$.
\end{theorem}

\begin{proof}
$E(n)$ "--- подгруппа, так как обе части сравнения $a^{(n-1)/2} \equiv \legendre{a}{n}$
мультипликативны по $a$. Покажем, что $E(n) \ne \Znx{n}$.

\emph{Случай 1: $p^2 \mid n$.} Эйлеров лжесвидетель "--- лжесвидетель Ферма, т.\,е.
$E(n) \subseteq L_F(n)$, а $L_F(n) \ne \Znx{n}$ по доказательству критерия Корсельта
(элемент $a = 1 + n/p$).

\emph{Случай 2: $n = p m$, $p \nmid m$, $m > 1$.} Пусть $g$ "--- квадратичный невычет по модулю $p$.
По китайской теореме найдётся $a$ с $a \equiv g \pmod p$ и $a \equiv 1 \pmod m$; оно взаимно
просто с $n$. Тогда $\legendre{a}{n} = \legendre{a}{p}\legendre{a}{m} = \legendre{g}{p}\legendre{1}{m} = -1$.
Если бы $a \in E(n)$, то $a^{(n-1)/2} \equiv -1 \pmod n$, и, приводя по модулю $m$,
$1 \equiv -1 \pmod m$ "--- невозможно при нечётном $m > 1$.
\end{proof}

\begin{algorithm}[тест Соловея"--~Штрассена]\label{l8:alg-ss}
\begin{algorithmic}[1]
\Require нечётное $n > 3$, число раундов $k$.
\Ensure «составное» (наверняка) или «вероятно простое».
\For{$i = 1, \ldots, k$}
  \State выбрать случайное $a \in \{2, \ldots, n - 2\}$
  \State $j \gets \legendre{a}{n}$ \Comment{алгоритм~\ref{l8:alg-jacobi}}
  \If{$j = 0$ или $a^{(n-1)/2} \not\equiv j \pmod n$} \Return «составное» \EndIf
\EndFor
\State \Return «вероятно простое»
\end{algorithmic}
\end{algorithm}

Раунд стоит $\bigO(\log^3 n)$ битовых операций, вероятность ошибки после $k$ раундов "---
не более $2^{-k}$. В отличие от теста Ферма, здесь нет аналога чисел Кармайкла: у любого
составного $n$ свидетелей Эйлера не меньше половины. Например, у числа Кармайкла $1729$
$a^{864} \equiv 1$ для \emph{всех} $a \in \Znx{1729}$, но $\legendre{a}{1729} = -1$ ровно для
половины из них ($648$ из $1296$), и все эти $a$ "--- свидетели Эйлера.

\taskexample{solovay-strassen-prime}

\paragraph{Сравнение тестов.} Можно доказать, что каждое сильно псевдопростое по основанию $a$
число является и эйлеровым псевдопростым по этому основанию (без доказательства,
см.~\cite[разд.~3.5]{crandall-pomerance}): свидетель Эйлера всегда является свидетелем Миллера.
Поэтому тест Миллера"--~Рабина не слабее теста Соловея"--~Штрассена, вдвое лучше по оценке
ошибки одного раунда ($1/4$ против $1/2$), не требует вычисления символа Якоби и при той же
стоимости раунда. Тест Соловея"--~Штрассена важен исторически: это первый
(1977) вероятностный алгоритм с доказанной оценкой ошибки, с которого началось изучение
вероятностных алгоритмов в теории сложности.

\begin{center}\small
\begin{tabular}{l|c|c|c}
  тест & условие & доля лжесвидетелей & «непробиваемые» $n$ \\ \hline
  Ферма & $a^{n-1} \equiv 1$ & $\le 1/2$ или $1$ & числа Кармайкла \\
  Соловей"--~Штрассен & $a^{(n-1)/2} \equiv \legendre{a}{n}$ & $\le 1/2$ & нет \\
  Миллер"--~Рабин & $a^d \equiv 1$ или $a^{2^i d} \equiv -1$ & $\le 1/4$ & нет
\end{tabular}
\end{center}

\section{Генерация случайных простых чисел}

Соберём всё вместе. Для генерации $b$-битного простого числа:

\begin{algorithm}[генерация простого числа]\label{l8:alg-gen}
\begin{algorithmic}[1]
\Require длина $b$, число раундов $k$.
\Ensure $b$-битное вероятно простое $n$.
\Repeat
  \State выбрать случайное нечётное $n \in [2^{b-1}, 2^b)$
  \If{$n$ делится на одно из небольших простых $3, 5, 7, \ldots, B$} \textbf{продолжить} \EndIf
  \Until{$n$ проходит $k$ раундов теста Миллера"--~Рабина}
\State \Return $n$
\end{algorithmic}
\end{algorithm}

По теореме~\ref{l8:pnt} в среднем требуется $\approx \frac12 \ln 2^b \approx 0{,}35\, b$
кандидатов; пробное деление на малые простые (дешёвое) отсеивает большую их часть, и тест
Миллера"--~Рабина применяется лишь к немногим. Все раунды, кроме последнего успешного
кандидата, обычно обрываются на первом же основании.

\checkyourself{Сколько в среднем кандидатов нужно проверить, чтобы найти $2048$-битное
простое число? Какую долю нечётных кандидатов отсеивает пробное деление на $3$, $5$ и $7$?}

\checkyourselfsection

\begin{furtherreading}
  \readitem[разд.~3.4, 3.9]{hps}{тест Миллера"--~Рабина, распределение простых чисел;
    квадратичные вычеты, символ Якоби}
  \readitem[гл.~1, разд.~3.4--3.5]{crandall-pomerance}{распределение простых чисел,
    псевдопростые, числа Кармайкла, сильные псевдопростые и теорема Монье"--~Рабина,
    тест Соловея"--~Штрассена}
  \readitem[гл.~3--5]{ireland-rosen}{теоремы Вильсона и Ферма, квадратичный закон взаимности}
  \readitem[гл.~22]{hardy-wright}{теорема Чебышёва и асимптотический закон}
  \readitem[гл.~10]{shoup}{тест Миллера"--~Рабина: доказательство оценки $1/4$}
  \readitem[разд.~4.2, 4.4]{hac}{вероятностные тесты простоты и генерация простых чисел}
  \readitem{msu-ent}{материалы кафедры теории чисел МГУ по элементарной теории чисел}
  \readitem{miller}{}
  \readitem{rabin}{}
  \readitem{solovay-strassen}{оригинальные статьи}
  \readitem{aks}{детерминированный полиномиальный тест простоты}
\end{furtherreading}
\end{lecture}
